\documentclass[11pt,a4paper]{amsart}

\usepackage[margin=1in]{geometry}
\usepackage[T1]{fontenc}
\usepackage[utf8]{inputenc}
\usepackage{lmodern}
\usepackage{microtype}
\usepackage{amsmath,amssymb,amsthm,mathtools}
\usepackage{enumitem}
\usepackage[colorlinks=true,linkcolor=blue,citecolor=blue,urlcolor=blue]{hyperref}
\usepackage[capitalize,noabbrev]{cleveref}

\numberwithin{equation}{section}

\newtheorem{theorem}{Theorem}[section]
\newtheorem{proposition}[theorem]{Proposition}
\newtheorem{lemma}[theorem]{Lemma}
\newtheorem{corollary}[theorem]{Corollary}
\newtheorem{definition}[theorem]{Definition}
\newtheorem{assumption}[theorem]{Assumption}
\theoremstyle{remark}
\newtheorem{remark}[theorem]{Remark}

\DeclareMathOperator{\RS}{RS}
\DeclareMathOperator{\Span}{Span}
\DeclareMathOperator{\Conf}{Conf}
\DeclareMathOperator{\supp}{supp}

\newcommand{\F}{\mathbb F}
\newcommand{\B}{\mathbb B}
\newcommand{\Z}{\mathbb Z}
\newcommand{\Ccal}{\mathcal C}
\newcommand{\Lcal}{\mathcal L}
\newcommand{\Scal}{\mathcal S}
\newcommand{\Wcal}{\mathcal W}
\newcommand{\OmegaC}{\Omega^{\circ}}
\newcommand{\Gord}{\Gamma^{\rm ord}}
\newcommand{\Gsid}{\Gamma^{\rm sid}}
\newcommand{\barN}{\overline N}
\newcommand{\eps}{\varepsilon}
\newcommand{\abs}[1]{\lvert #1\rvert}
\newcommand{\defeq}{\coloneqq}

\title{Asymptotic RS MCA}
\date{}

\begin{document}

\begin{abstract}
We give a compact conditional proof of the asymptotic Reed--Solomon mutual-correlated-agreement frontier from the paid ledger developed in the v13 raw and Grande Finale manuscripts.  The proof reduces the remaining primitive prefix problem to a fixed-density Boolean slice mapped by Vandermonde moments.  After the image-normalized Fourier/Sidon cell is paid, any positive logarithmic moment excess gives a large high-energy Boolean fiber.  Balog--Szemeredi--Gowers and quasicube difference growth forbid such a fiber, so primitive Q holds with only \(\exp(o(n))\) loss.  Combined with the stated add-back and lower-side inputs, every smooth row sequence covered by the conditional closed ledger has threshold \(1-\rho-g^*(\rho,\log_2|\B|)+o(1)\).  The closed-ledger package is stated explicitly, with the individual cell payments cited to the existing ledgers and with the C9 normalization exposed as an assumption.
\end{abstract}

\maketitle

\section{Statement}

Let \(D\subseteq \F\) be an evaluation set of size \(n\), and let
\[
        C=\RS_{\F}(D,k)=\{(f(t))_{t\in D}:\deg f<k\}.
\]
For \(a\in\{0,\ldots,n\}\), a finite slope \(\gamma\in\F\) is support-wise MCA-bad for a pair \(r_1,r_2\in\F^D\) at agreement \(a\) if some set \(S\subseteq D\), \(|S|\ge a\), and some \(h\in\F[X]\), \(\deg h<k\), satisfy \(r_1+\gamma r_2=h\) on \(S\), while \(r_1,r_2\) are not simultaneously explained on the same \(S\) by two degree-\(<k\) polynomials.  Put
\[
        B_C^{\rm MCA}(a)
        =\max_{r_1,r_2}\#\{\gamma\in\F:\gamma\text{ is MCA-bad at agreement }a\}.
\]
Then \(\eps_{\rm mca}(C,1-a/n)=B_C^{\rm MCA}(a)/q_{\rm line}\), where \(q_{\rm line}\) is the line-slope field used by the MCA sampler.  This integer staircase convention is the one used in the v13 raw ledger \cite{Cho26CapV13} and in the Grande Finale reduction \cite{Cho26Grande}.

We consider smooth multiplicative or circle-type RS row sequences \(C_n=\RS_{\F_n}(D_n,k_n)\) with \(|D_n|=n\), \(k_n/n\to\rho\in(0,1)\), and a base/generated alphabet \(\B_n\) of bit-size \(\beta_n=\log_2|\B_n|\to\beta\).  For an agreement \(a_n=k_n+1+w_n\), define the prefix average
\[
        \barN_{n,a_n}=\binom n{a_n}|\B_n|^{-w_n}.
\]
The entropy--subfield envelope is
\[
        g^*(\rho,\beta)=\sup\{g\in[0,1-\rho]:H_2(\rho+g)\ge \beta g\},
\]
where \(H_2(x)=-x\log_2x-(1-x)\log_2(1-x)\).  The target threshold is \(\delta_C^*(\eps)=\sup\{\delta:\eps_{\rm mca}(C,\delta)\le\eps\}\).

\begin{definition}[Cells and first-match payment]\label{def:cells}
A \emph{cell} is a constructible subfamily of the MCA witness incidence space, or a bounded-complexity family of such subfamilies, stable under extension of the ground field.  Cells are ordered.  A slope is assigned to the first cell containing one of its witnesses.  The cell is \emph{paid} at agreement \(a_n\) if the number of slopes assigned to it is at most its printed budget, and a family is paid if the sum of its first-match budgets is \(\exp(o(n))\barN_{n,a_n}\).
\end{definition}

\begin{lemma}[First-match disjointization]\label{lem:first-match}
If ordered cells \(\Ccal_1,\ldots,\Ccal_J\) cover a set of bad slopes and the \(j\)-th first-match cell has budget \(U_j\), then the covered slopes contribute at most \(\sum_jU_j\) to \(B_C^{\rm MCA}(a)\).
\end{lemma}

\begin{proof}
For a fixed received line, assign each bad slope to the least-indexed cell containing a witness.  These assigned slope classes are disjoint, and the \(j\)-th class is contained in the projection of \(\Ccal_j\) after earlier cells have been removed.  Summing the budgets and then maximizing over received lines proves the claim.
\end{proof}

\begin{definition}[Closed paid ledger]\label{def:closed-ledger}
A smooth RS row sequence is \emph{closed-ledger} at agreements \(a_n\) if the following first-match cells cover all MCA-bad witnesses except the primitive leaves of \cref{def:primitive-leaf}, and if their total contribution is \(\exp(o(n))\barN_{n,a_n}\).
\begin{enumerate}[label=\textup{(C\arabic*)},leftmargin=2.7em]
    \item \emph{Quotient-pullback cells.}  The support or locator descends along a nontrivial finite map \(D\to D'\).  These cells include quotient-remainder locators, divisor-block support ledgers, quotient-image lcm ledgers, and quotient pullback recursion.
    \item \emph{Chebyshev and dihedral cells.}  On multiplicative or circle rows, the support is invariant under the dihedral action, or the locator factors through a Chebyshev or circle quotient such as \(z+z^{-1}\).
    \item \emph{Planted-block cells.}  A positive-density block is forced into the support or common zero set.  The payment is the entropy cost of planting that block, together with the bounded quotient profile.
    \item \emph{Tangent and deep-center cells.}  The received line lies in a tangent, secant, or high-contact stratum of the RS code variety, including the high-agreement tangent cell.
    \item \emph{Extension and descent cells.}  The witness either descends to a smaller field or is controlled by a bounded extension-line or extension-pole parameter.
    \item \emph{Differential-locator and regular-rank cells.}  The locator map has a rank drop, a ramification defect, or an overdetermined Hankel minor.  These cells include regular closed-range Hankel packing, canonical rank-drop ledgers, and curve rank-drop ledgers.
    \item \emph{Saturation and effective-image-collapse cells.}  The apparent support count is larger than the actual line-ray or codeword-ray image.  The payment uses exact saturation identities and exact residual image ledgers.
    \item \emph{Balanced-core and split-pencil cells.}  The residual witness lies in a balanced-core split-pencil chart.  One-parameter primitive locator pencils are paid by the moving-root theorem; deficiency-one and identically split top charts are tangent-borne or eliminated.
    \item \emph{Fourier/Sidon cells.}  In a primitive prefix leaf, an image-normalized heavy fiber with exponentially small additive closure is cut as a Fourier/Sidon cell.  The moment is the \(L=|\operatorname{im}\Phi|\) average from \cref{def:primitive-leaf}, not an ambient \(|K|^R\) average: major arcs are routed to \textup{(C1)--(C8)}, and minor arcs are paid by a Fourier/Sidon theorem stated at image scale.
\end{enumerate}
The remaining first-match complement is a disjoint union of primitive prefix leaves.  Every such leaf has the Boolean fixed-density Vandermonde form of \cref{def:primitive-leaf}.
\end{definition}

\begin{theorem}[Conditional closed ledger package]\label{thm:closed-ledger-package}
For the smooth multiplicative and circle-type RS row sequences treated in the v13 raw ledger and in the Grande Finale notes, assume the cited \textup{(C1)--(C8)} payments hold uniformly in the asymptotic frontier window, and assume the image-normalized C9 input of \cref{ass:image-normalized-sidon-input} for every surviving primitive leaf.  Then the cells \textup{(C1)--(C9)} form a closed paid ledger in that window.
\end{theorem}

\begin{proof}
We record the existing ledger package rather than reprove each algebro-geometric count.

For \textup{(C1)}, the quotient-remainder lower and support ledgers, divisor-union support ledger, exact finite-parameter quotient-image ledger, affine/projective lcm ledgers, quotient pullback recursion, and bounded quotient census are proved in \cite[quotient-remainder, quotient support, quotient-image, and bounded-census ledgers]{Cho26CapV13}.  The corresponding first-match convention and quotient status cell are isolated in \cite{Cho26Grande}.

For \textup{(C2)}, exact Chebyshev fibers on circle twin-coset domains, circle quotient transport, and circle deployed ledgers are proved in \cite[Chebyshev and circle sections]{Cho26CapV13}.  Dihedral and Chebyshev cells are quotient cells with extra monodromy and are therefore paid by the same downstairs-support and image ledgers.

For \textup{(C3)}, planted quotient profiles, planted quotient-core lower counts, planted list-budget windows, and bounded active quotient order are established in \cite[planted quotient profile and bounded quotient census]{Cho26CapV13}.  These supply the entropy payment for positive-density forced blocks.

For \textup{(C4)}, the exact high-agreement tangent cell and tangent-borne split-kernel sections are proved in \cite[high-agreement tangent exact cell and split-kernel tangent-borne theorem]{Cho26CapV13} and refined in \cite[exact high-agreement tangent cell]{Cho26Grande}.

For \textup{(C5)}, the extension-pole deep-list and quotient-remainder floors, extension-line dimension-degree ledger, subfield confinement, and genuinely extension-valued line accounting are proved in \cite[extension-pole and extension-line ledgers]{Cho26CapV13}; the full-extension cell target and extension-valued rank-one floor are recorded in \cite{Cho26Grande}.

For \textup{(C6)}, regular Hankel minors, canonical exact-agreement rank-drop ledgers, closed-ball rank-drop lcm ledgers, curve rank-drop ledgers, arbitrary-residual image ledgers, and the scanner residual ledger are proved in \cite[Hankel, rank-drop, curve, residual-image, and scanner ledgers]{Cho26CapV13}.  These are exactly the differential-locator and regular-rank payments.

For \textup{(C7)}, the exact saturation identity and line-ray saturation identity are proved in \cite[saturation and line-ray saturation identities]{Cho26Grande}; exact arbitrary-residual image and closed-ball residual image ledgers appear in \cite[residual image ledgers]{Cho26CapV13}.  Thus raw support overcounts are replaced by actual image counts.

For \textup{(C8)}, the interpolation-lattice split-pencil reduction, near-rational payment, deficiency-one split-test eliminant, split annihilator dictionary, tangent-borne split charts, boundary-profile identification with Q, base-field split-pencil floor, one-parameter moving-root theorem, and primitive one-pencil corollary are proved in \cite[split-pencil and BC sections]{Cho26Grande}, with the balanced-core and split-pencil reductions also developed in \cite[rank-one and split-pencil reductions]{Cho26CapV13}.

For \textup{(C9)}, Fourier-flat leaves satisfy asymptotic Q and large-characteristic Fourier examples are proved in \cite[Fourier-flat Q section]{Cho26Grande}.  The Sidon/Fourier operational cut, the split between Sidon-heavy and high-energy fibers, and the fact that major arcs route to algebraic cells while minor arcs are Fourier-flat are the moduli-ledger results of \cite{Cho26ModuliFinal,Cho26ModuliSelf}.  The C9 input consumed here is the image-normalized assumption in \cref{ass:image-normalized-sidon-input}.  Ambient max-fiber information may enter through \cref{lem:ambient-image-max}, and ambient moment information may enter through \cref{lem:moment-normalization}; no unprinted bridge \(A/L=\exp(o(N))\) is used.  The image-normalized C9 input is used below in \cref{def:sidon-paid,prop:energy-extract,thm:primitive-q}.

Finally, first-match disjointization is \cref{lem:first-match} here and is also proved in the Grande Finale and moduli-ledger notes.  The listed cells exhaust the nonprimitive branches of the v13 first-match atlas; after removing them, the remaining leaf is the primitive Boolean prefix slice.  Summing the printed budgets gives \(\exp(o(n))\barN_{n,a_n}\), which is the closed paid ledger.
\end{proof}

Thus the row-specific geometry is packaged as an imported conditional ledger.  The internal work below starts after those cell inputs, the image-normalized C9 assumption, the add-back profile input in \cref{lem:addback}, and the lower-side input in the proof of \cref{thm:frontier}.

\begin{theorem}[Asymptotic RS--MCA frontier]\label{thm:frontier}
Let \(C_n=\RS_{\F_n}(D_n,k_n)\) be a smooth RS row sequence as above.  Assume the conditional closed paid ledger of \cref{thm:closed-ledger-package} holds uniformly in every \(o(n)\)-window around the crossing \(a_n=(\rho+g^*)n+o(n)\).  Assume also the add-back profile input in \cref{lem:addback} and a lower-side identity-prefix input with subexponential collision loss, or an equivalent collision-free identity-prefix floor.  Then, for every target sequence with \(\log_2(1/\eps_n)=O(n)\) and fixed challenge normalization,
\[
        \delta_{C_n}^*(\eps_n)
        =1-\rho-g^*(\rho,\beta)+o(1),
\]
with the evident replacement of \(\beta\) by \(\lim\log_2|\B_n|\) when the base alphabet varies.
\end{theorem}

\begin{remark}[Current audit status]\label{rem:current-audit-status}
The normalization repair in this version is limited to B1: C9 is consumed at image scale, and ambient/image moment transfers are printed explicitly.  The note does not supply the missing C9 moduli/Fourier--Sidon source theorem, the B3 window-uniformity theorem, the add-back profile decomposition, or the lower-side pole-reservoir collision-loss proof.
\end{remark}

Subject to the imported cell inputs, the image-normalized C9 assumption, the add-back profile input, and the lower-side input stated above, the rest of the paper proves \cref{thm:frontier}.  The proof is otherwise self-contained after the standard external additive-combinatorics facts stated in \cref{sec:bsg}.

\section{Primitive leaves}

\begin{definition}[Primitive prefix leaf]\label{def:primitive-leaf}
A primitive leaf consists of a finite set \(T\), \(|T|=N\), an integer \(m\) with \(m/N\) bounded away from \(0\) and \(1\), a field \(K\), nonzero weights \(\rho(t)\in K^\times\), and columns
\[
        v_t=\rho(t)(1,t,\ldots,t^{R-1})\in K^R,\qquad t\in T.
\]
The active support family is a first-match residual slice
\[
        \OmegaC\subseteq \{x\in\{0,1\}^{T}:\sum_{t\in T}x_t=m\},
\]
and the prefix map is \(\Phi(x)=\sum_{t\in T}x_tv_t\).  Write \(\Scal=\operatorname{im}\Phi\), \(L=|\Scal|\), \(M=|\OmegaC|\), \(\barN=M/L\), and \(F_s=\OmegaC\cap\Phi^{-1}(s)\).  The leaf is at the frontier if \(\log\barN=o(N)\).
\end{definition}

For \(q=q_N\to\infty\), define
\[
        \Gord_q=L^{-1}\sum_{s\in\Scal}\left(\frac{|F_s|}{\barN}\right)^q .
\]

\begin{lemma}[Moment-max equivalence]\label{lem:moment-max}
If \(q\to\infty\) and \(\log L=o(Nq)\), then
\[
        L^{-1}\left(\max_s\frac{|F_s|}{\barN}\right)^q
        \le \Gord_q
        \le
        \left(\max_s\frac{|F_s|}{\barN}\right)^q .
\]
Hence \(\Gord_q\le \exp(o(Nq))\) is equivalent to \(\max_s|F_s|\le\exp(o(N))\barN\).
\end{lemma}

\begin{proof}
The upper bound replaces every summand by the maximum; the lower bound keeps a maximum summand.  Taking \(q\)-th roots and using \(\log L=o(Nq)\) gives the final assertion.
\end{proof}

\begin{lemma}[Max-fiber ambient-to-image transfer]\label{lem:ambient-image-max}
Let \(\Phi:\OmegaC\to G\) map into a finite abelian ambient group \(G\) of size \(A\).  Put \(\Scal=\operatorname{im}\Phi\), \(L=|\Scal|\), \(M=|\OmegaC|\), and \(F_s=\OmegaC\cap\Phi^{-1}(s)\).  If
\[
        \max_{y\in G}|\OmegaC\cap\Phi^{-1}(y)|
        \le \exp(o(N))\frac{M}{A},
\]
then
\[
        \max_{s\in\Scal}|F_s|
        \le \exp(o(N))\frac{M}{L}.
\]
\end{lemma}

\begin{proof}
Every image fiber is an ambient fiber, so
\[
        \max_{s\in\Scal}|F_s|
        \le \max_{y\in G}|\OmegaC\cap\Phi^{-1}(y)|
        \le \exp(o(N))\frac{M}{A}.
\]
Since \(L\le A\), one has \(M/A\le M/L\), proving the claimed image-normalized bound.
\end{proof}

\begin{lemma}[Moment normalization identity]\label{lem:moment-normalization}
Let \(\Phi:\OmegaC\to G\) map into a finite abelian ambient group \(G\) of size \(A\).  Put \(\Scal=\operatorname{im}\Phi\), \(L=|\Scal|\), \(M=|\OmegaC|\), \(\barN=M/L\), and \(F_s=\OmegaC\cap\Phi^{-1}(s)\).  For \(q\ge1\), define
\[
        \Gamma_{\rm img}(q)
        =L^{-1}\sum_{s\in\Scal}\left(\frac{|F_s|}{M/L}\right)^q
\]
and
\[
        \Gamma_{\rm amb}(q)
        =A^{-1}\sum_{y\in G}
        \left(\frac{|\OmegaC\cap\Phi^{-1}(y)|}{M/A}\right)^q .
\]
Then
\[
        \Gamma_{\rm amb}(q)
        =\left(\frac{A}{L}\right)^{q-1}\Gamma_{\rm img}(q).
\]
Consequently,
\[
        \Gamma_{\rm img}(q)
        =\left(\frac{L}{A}\right)^{q-1}\Gamma_{\rm amb}(q)
        \le \Gamma_{\rm amb}(q)
        \qquad(q\ge1),
\]
so an ambient upper bound is safe to use as an image upper bound.  The reverse direction is unsafe without a printed bound on \(A/L\), and the Fourier/Sidon payment consumed below is therefore stated directly at image scale.
\end{lemma}

\begin{proof}
The summands with \(y\notin\Scal\) vanish.  Since \(M/A=(L/A)(M/L)=(L/A)\barN\),
\[
\begin{aligned}
        \Gamma_{\rm amb}(q)
        &=A^{-1}\sum_{s\in\Scal}
        \left(\frac{|F_s|}{M/A}\right)^q \\
        &=A^{-1}\left(\frac{A}{L}\right)^q
        \sum_{s\in\Scal}\left(\frac{|F_s|}{\barN}\right)^q \\
        &=\left(\frac{A}{L}\right)^{q-1}
        L^{-1}\sum_{s\in\Scal}\left(\frac{|F_s|}{\barN}\right)^q .
\end{aligned}
\]
This is the displayed identity.
\end{proof}

\section{Sidon cut}

For a finite \(F\) in an abelian group, let
\[
        E(F)=\#\{(a,b,c,d)\in F^4:a-b=c-d\},\qquad
        \Delta(F)=E(F)/|F|^3.
\]
Thus \(\Delta(F)\) is the probability that \(a-b+c\in F\) for independent uniform \(a,b,c\in F\).

For \(\sigma>0\), define the Sidon-heavy moment of a primitive leaf by
\[
        \Gsid_{q,\sigma}
        =L^{-1}\sum_{\Delta(F_s)\le e^{-\sigma N}}
        \left(\frac{|F_s|}{\barN}\right)^q .
\]

\begin{definition}[Paid Fourier/Sidon cell]\label{def:sidon-paid}
The Fourier/Sidon cell is paid if, for every fixed \(\sigma>0\) and every logarithmic \(q\), one has \(\Gsid_{q,\sigma}\le\exp(o(Nq))\).  Equivalently, no positive-rate heavy fiber with exponentially small additive closure remains in the primitive residual.  This is an image-normalized statement: \(L=|\operatorname{im}\Phi|\), \(\barN=M/L\), and the average is over \(\Scal\), not over an ambient group.
\end{definition}

\begin{assumption}[Image-normalized C9 Fourier/Sidon input]\label{ass:image-normalized-sidon-input}
Let \((T,m,K,\rho(t),v_t,\OmegaC,\Phi)\) be a primitive prefix leaf surviving cells \textup{(C1)--(C8)} in the asymptotic frontier window.  Write
\[
        \Scal=\operatorname{im}\Phi,\qquad
        L=|\Scal|,\qquad
        M=|\OmegaC|,\qquad
        \barN=M/L,\qquad
        F_s=\OmegaC\cap\Phi^{-1}(s).
\]
The C9 moduli ledger input used by this note is the following image-scale payment: for every fixed \(\sigma>0\) and every logarithmic \(q=q_N\to\infty\),
\[
        \Gsid_{q,\sigma}
        =L^{-1}\sum_{\Delta(F_s)\le \exp(-\sigma N)}
        \left(\frac{|F_s|}{\barN}\right)^q
        \le \exp(o(Nq)).
\]
Equivalently, the Fourier/Sidon cell is paid in the image normalization used by \cref{def:primitive-leaf}.
\end{assumption}

This is the C9 Fourier/Sidon input stated in the leaf normalization of \cref{def:primitive-leaf}.  The present note does not prove the missing moduli ledger; it records the exact normalization consumed downstream.  A proof of C9 may provide the displayed inequality directly, or may first prove a stronger image-scale minor-arc estimate after major arcs have been routed to \textup{(C1)--(C8)}.  Ambient moment estimates can be consumed here only through \cref{lem:ambient-image-max,lem:moment-normalization}.

In practice this cell is paid by Fourier flatness after major arcs have been routed to algebraic cells.  The output needed here is an image-scale bound such as \(L\mu(s)\le\exp(o(N))\) for all \(s\in\Scal\), where \(\mu=\Phi_*\operatorname{Unif}(\OmegaC)\), or directly the bound in \cref{ass:image-normalized-sidon-input}.  An ambient Fourier estimate does not imply this image-scale statement without one of the printed transfers above.  In RS moment-curve leaves, large Fourier coefficients are algebraic major arcs: quotient, descent, dihedral, ramification, or effective-image-collapse phenomena.  These are exactly the earlier paid cells in \cref{def:closed-ledger}.

\begin{proposition}[Energy extraction after the Sidon cut]\label{prop:energy-extract}
Assume the Fourier/Sidon cell is paid.  If \(\Gord_q\ge e^{\eta Nq}\) for some fixed \(\eta>0\), then some level \(s=s_N\) satisfies \(|F_s|\ge e^{\eta N-o(N)}\barN\) and \(E(F_s)\ge |F_s|^3e^{-o(N)}\).
\end{proposition}

\begin{proof}
Fix \(\sigma>0\).  Split \(\Gord_q\) into terms with \(\Delta(F_s)\le e^{-\sigma N}\) and the complementary terms.  The first part is \(\Gsid_{q,\sigma}=\exp(o(Nq))\), so a complementary term has \((|F_s|/\barN)^q\ge e^{\eta Nq-o(Nq)}\).  Hence \(|F_s|\ge e^{\eta N-o(N)}\barN\) and \(\Delta(F_s)>e^{-\sigma N}\).  Letting \(\sigma\downarrow0\) slowly along the sequence gives \(E(F_s)\ge |F_s|^3e^{-o(N)}\).
\end{proof}

\section{Boolean additive combinatorics}\label{sec:bsg}

We use the energy form of Balog--Szemeredi--Gowers \cite{BalogSzemeredi,Gowers1998}.  Its precise polynomial exponents are irrelevant here.

\begin{theorem}[Balog--Szemeredi--Gowers]\label{thm:bsg}
If \(A\) is a finite subset of an abelian group and \(E(A)\ge |A|^3/K\), then some \(A'\subseteq A\) satisfies \(|A'|\ge K^{-C}|A|\) and \(|A'-A'|\le K^C|A'|\), where \(C\) is absolute.
\end{theorem}

We also use the quasicube sumset theorem of Matolcsi--Ruzsa--Shakan--Zhelezov \cite{MRSZ2020} in the sharpened form of Green--Matolcsi--Ruzsa--Shakan--Zhelezov \cite{GMRSZ2020}: if \(U\subseteq\{0,1\}^d\), then \(|P+Q+U|\ge |P|^{1/2}|Q|^{1/2}|U|\) for finite nonempty \(P,Q\subseteq\Z^d\).

\begin{theorem}[Quasicube difference growth]\label{thm:quasicube}
Every finite \(A\subseteq\{0,1\}^{N}\) satisfies \(|A-A|\ge |A|^{3/2}\).
\end{theorem}

\begin{proof}
Apply the quasicube theorem with \(U=A\), \(P=-A\), and \(Q=\{0\}\).
\end{proof}

\begin{proposition}[No large high-energy Boolean fiber]\label{prop:no-high-energy}
There is no sequence \(F_N\subseteq\{0,1\}^{N}\) with \(|F_N|\ge e^{cN-o(N)}\) and \(E(F_N)\ge |F_N|^3e^{-o(N)}\) for fixed \(c>0\).
\end{proposition}

\begin{proof}
By \cref{thm:bsg} with \(K=e^{o(N)}\), a subset \(F'_N\subseteq F_N\) has \(|F'_N|\ge e^{cN-o(N)}\) and \(|F'_N-F'_N|\le e^{o(N)}|F'_N|\).  But \cref{thm:quasicube} gives \(|F'_N-F'_N|\ge |F'_N|^{3/2}\), so \(|F'_N|^{1/2}\le e^{o(N)}\), a contradiction.
\end{proof}

\begin{theorem}[Primitive Q after the Sidon cut]\label{thm:primitive-q}
In every frontier primitive leaf with paid Fourier/Sidon cell, \(\max_s|F_s|\le\exp(o(N))\barN\).
\end{theorem}

\begin{proof}
If not, \cref{lem:moment-max} gives \(\Gord_q\ge e^{\eta Nq}\) for some fixed \(\eta>0\) along a subsequence and some logarithmic \(q\).  \Cref{prop:energy-extract} then gives a fiber \(F_s\subseteq\{0,1\}^{T}\) with \(|F_s|\ge e^{\eta N-o(N)}\barN=e^{cN-o(N)}\) and \(E(F_s)\ge |F_s|^3e^{-o(N)}\), because \(\log\barN=o(N)\).  This contradicts \cref{prop:no-high-energy}.
\end{proof}

\section{MCA compiler}

\begin{lemma}[Conditional first-match add-back]\label{lem:addback}
Assume the paid primitive leaves admit a first-match profile decomposition into \(\exp(o(n))\) profiles, and that the normalized primitive leaf averages in those profiles sum to at most \(\exp(o(n))\barN_{n,a_n}\).  Then, for a closed-ledger row, the full first-match prefix count after paid cells satisfies \(\max_sN(s)\le\exp(o(n))\barN_{n,a_n}\).
\end{lemma}

\begin{proof}
Paid cells contribute \(\exp(o(n))\barN_{n,a_n}\) by closure.  Primitive leaves satisfy \cref{thm:primitive-q}.  The stated profile-decomposition input bounds the sum of their normalized averages, and summing over the subexponential profile family preserves the same scale.  The present note records this add-back dependency rather than proving the profile decomposition.
\end{proof}

\begin{lemma}[Q pays shift pairs]\label{lem:q-sp}
If \(\sum_sN(s)=M\), \(\barN=M/L\), and \(\max_sN(s)\le \kappa\barN\), then \(M^{-1}\sum_sN(s)^2\le\kappa\barN\).
\end{lemma}

\begin{proof}
\(\sum_sN(s)^2\le(\max_sN(s))\sum_sN(s)\le\kappa\barN M\).
\end{proof}

\begin{theorem}[Conditional closed-ledger upper bound]\label{thm:upper}
If a smooth RS row sequence is closed-ledger at \(a_n\) and satisfies the add-back input in \cref{lem:addback}, then
\[
        B_{C_n}^{\rm MCA}(a_n)
        \le \exp(o(n))\binom n{a_n}|\B_n|^{-(a_n-k_n-1)} .
\]
\end{theorem}

\begin{proof}
The nonprimitive cells are paid by definition.  The primitive prefix leaves obey Q by \cref{thm:primitive-q} and add back by \cref{lem:addback}.  The shift-pair second-moment ledger is paid by \cref{lem:q-sp}.  First-match disjointization sums these disjoint slope classes, giving the displayed numerator bound.
\end{proof}

\section{Lower side and frontier}

We recall the identity-prefix pole construction in the form needed here.  Let \(m=k+1+w\), and for \(S\subseteq D\), \(|S|=m\), write \(\ell_S(X)=\prod_{t\in S}(X-t)\).  Pigeonholing the first \(w\) locator coefficients over \(\B^w\) gives a prefix vector \(z\) with at least \(\binom nm|\B|^{-w}\) supports.  Choose a pole \(\alpha\in\F\setminus D\), and set \(U_z(X)=X^m+z_1X^{m-1}+\cdots+z_wX^{m-w}\).  The line
\[
        f_\alpha(x)=\frac{U_z(x)}{x-\alpha},
        \qquad
        g_\alpha(x)=\frac{-1}{x-\alpha},
        \qquad x\in D,
\]
has the property that \(f_\alpha+\zeta g_\alpha\) is code-explained on \(S\) exactly when \(\ell_S(\alpha)=U_z(\alpha)-\zeta\).  The word \(g_\alpha\) is not explained by a degree-\(<k\) polynomial on \(k+1\) positions, so the corresponding slopes are MCA-bad.  The lower side consumed by \cref{thm:frontier} assumes either a subexponential collision loss for evaluating \(\ell_S(\alpha)\) in the pole-reservoir regime, or a replacement by a collision-free identity-prefix floor.  Under that lower-side input, this gives
\[
        B_C^{\rm MCA}(m)\ge \exp(-o(n))\binom nm|\B|^{-w}.
\]
This is the lower construction used in \cite{Cho26CapV13,Cho26Grande}.

\begin{proof}[Proof of \cref{thm:frontier}]
Take \(a_n=(\rho+g)n+o(n)\).  Stirling gives
\[
        \log_2\barN_{n,a_n}
        = n\bigl(H_2(\rho+g)-\beta g\bigr)+o(n).
\]
If \(g>g^*(\rho,\beta)\) by a fixed amount, this exponent is negative by \(\Omega(n)\), so \cref{thm:upper} makes \(B_{C_n}^{\rm MCA}(a_n)\) smaller than any target with the prescribed normalization after moving \(o(n)\) agreements to absorb the \(\exp(o(n))\) ledger overhead.  Thus the safe radius reaches \(1-\rho-g^*(\rho,\beta)-o(1)\).  If \(g<g^*(\rho,\beta)\) by a fixed amount, the identity-prefix construction gives exponentially many bad slopes, so the target fails at radius \(1-\rho-g+o(1)\).  Letting \(g\) approach \(g^*\) from both sides proves the claimed equality.
\end{proof}

\begin{remark}[Finite adjacent rows]
This is an asymptotic theorem.  It intentionally absorbs \(\exp(o(n))\) losses.  The finite deployed adjacent pairs in \cite{Cho26CapV13,Cho26Grande} require the same first-match ledger with printed constants for every paid cell, Fourier/Sidon tail, and add-back step; those constants are not supplied by the asymptotic argument alone.
\end{remark}

\begin{thebibliography}{GMRSZ20}

\bibitem[BS94]{BalogSzemeredi}
A. Balog and E. Szemeredi,
\newblock A statistical theorem of set addition,
\newblock \emph{Combinatorica} 14 (1994), 263--268.

\bibitem[Cho26a]{Cho26CapV13}
P. Chojecki,
\newblock \emph{CAP25 v13 raw ledger for smooth Reed--Solomon MCA},
\newblock experimental manuscript, 2026.

\bibitem[Cho26b]{Cho26Grande}
P. Chojecki,
\newblock \emph{Grande Finale: Q, BC, SP, and the final RS--MCA ledger},
\newblock experimental manuscript, 2026.

\bibitem[Cho26c]{Cho26ModuliSelf}
P. Chojecki,
\newblock \emph{The Moduli Ledger for the Asymptotic Reed--Solomon MCA Frontier},
\newblock experimental manuscript, 2026.

\bibitem[Cho26d]{Cho26ModuliFinal}
P. Chojecki,
\newblock \emph{A Moduli Ledger for RS--MCA},
\newblock experimental manuscript, 2026.

\bibitem[GMRSZ20]{GMRSZ2020}
B. Green, D. Matolcsi, I. Z. Ruzsa, G. Shakan, and D. Zhelezov,
\newblock A weighted Prekopa--Leindler inequality and sumsets with quasicubes,
\newblock arXiv:2003.04077, 2020.

\bibitem[Gow98]{Gowers1998}
W. T. Gowers,
\newblock A new proof of Szemeredi's theorem for arithmetic progressions of length four,
\newblock \emph{Geom. Funct. Anal.} 8 (1998), 529--551.

\bibitem[MRSZ20]{MRSZ2020}
D. Matolcsi, I. Z. Ruzsa, G. Shakan, and D. Zhelezov,
\newblock An analytic approach to cardinalities of sumsets,
\newblock arXiv:2003.04075, 2020.

\end{thebibliography}

\end{document}
